%% kapa: exam=1 round=22 year=2011 period=1 form=A subject=civil_law subject_ko=민법 seq=31 number=31
%% source: Q-Net 감정평가사 기출문제 (cst00309) artlSeq=5036790 file=extracted/1008/5036790/2011년 제22회 감정평가사 1차시험 기출문제/1교시.hwp
%% conversion: kordoc 4.24.51 -> tools/gen1.py (auto); review=proofread-sonnet; flags=none
\begin{question}{31}
甲은 법령상 토지거래허가구역 내의 자신의 토지에 대해 허가를 받지 않은 채 乙과 매매계약을 체결하고 계약금을 지급받았다. 다음 설명으로 옳지 않은 것은? (다툼이 있으면 판례에 의함)
\begin{choices}
  \item 甲이 乙에게 소유권이전등기를 해 주지 않는 경우라도 토지거래허가를 받기 전에는 乙은 채무불이행을 이유로 계약을 해제할 수 없다.
  \item 甲은 토지거래허가를 받지 않은 상태에서는 특별한 사정이 없는 한 乙에게 계약금의 배액을 상환하고 매매계약을 해제할 수 없다.
  \item 甲과 乙 사이의 매매계약은 허가받기 전의 상태에서는 무효이므로, 乙이 허가받을 것을 조건으로 소유권이전등기청구권을 피보전권리로 한 부동산처분금지가처분 신청을 하는 것은 허용되지 않는다.
  \item 甲이 협력의무를 이행하지 않음으로써 乙이 손해를 입은 경우, 乙은 甲에게 손해배상을 청구할 수 있을 뿐, 협력의무 불이행을 이유로 매매계약 자체를 해제할 수 없다.
  \item 만약 허가요건을 갖추지 못한 乙이 허가요건을 갖춘 사람의 명의를 도용하여 매매계약서에 그를 매수인으로 기재한 경우, 그 매매계약은 처음 체결된 때부터 확정적으로 무효이다.
\end{choices}
\end{question}
\answer{2}
